Classical unraveling further predicts that, once disclosure is mandatory and verifiable, the share of non-disclosing listings should converge rapidly toward zero. The French data reject this prediction in both segments, though with different magnitudes. Figure~\ref{fig:fact2_missing} plots the monthly share of listings without a certified EPC from October 2018 through July 2024. For apartments (panel~(a)), the share averages $35.9\%$ before the reform and $33.8\%$ afterward; over the final twelve months of the sample, it stabilizes at $28.9\%$. For houses (panel~(b)), the corresponding numbers are $24.6\%$ pre-reform, $14.4\%$ post-reform, and a final-twelve-month mean of $9.7\%$. Houses thus display better baseline audit coverage and faster post-reform convergence, but in both segments the share remains structurally bounded away from zero more than three years after disclosure became legally required. Almost three-tenths of apartment listings and one-tenth of house listings still come to market without a certified audit, indicating that a substantial fraction of sellers cannot disclose, because they do not possess a certified audit at the time of listing.
